% ============================================================
%  TEMPLATE — Apostila PratiqueJá
%  Engine: XeLaTeX
%  Como usar: copie este arquivo, renomeie e edite apenas
%             o conteúdo dentro de \begin{multicols}.
% ============================================================
\documentclass[12pt]{article}
\usepackage{pratiqueja}

% ── Defina o assunto que aparece no rodapé central ────────────
\setsubject{Nome do Assunto}

\begin{document}
\pagenumbering{arabic}
\setstretch{1.2}
\color{bodytext}

% ── Cabeçalho: {Título}{Subtítulo}{Área · Nível} ─────────────
\heroheader{Título do Assunto}%
           {Subtítulo ou descrição breve}%
           {Matemática · Intermediário}

\setlength{\columnsep}{18pt}
\setlength{\columnseprule}{0.5pt}
\def\columnseprulecolor{\color{exborder}}

\begin{multicols}{2}

% ── Seção: \TW{badge}{cor}{título}{subtítulo} ─────────────────
% Cores disponíveis: darkblue | badgepurple | badgegreen
%
% \TW{Def}{badgepurple}{Definição}{Nome do conceito}
% \regra{Texto da regra ou definição.}
%
% \TW{1}{darkblue}{Primeiro tópico}{Descrição curta}
% \regra{Regra do tópico.}
% \exemp{%
%   Exemplo 1 → resultado \ok\\
%   Exemplo 2 → resultado \nok%
% }
% \obs{Observação complementar.}
%
% \formula{$f(x) = ...$}
%
% \SH{Rótulo discreto de sub-seção}
%
% \chip{A}\chip{B}\chip{C}   ← chips de itens

% ─────────────────────────────────────────────────────────────
% ESCREVA O CONTEÚDO ABAIXO DESTA LINHA
% ─────────────────────────────────────────────────────────────



\end{multicols}
\end{document}
